\documentclass[10pt,a4paper]{amsart}
\usepackage[margin=19mm]{geometry}
\usepackage{amsmath,amssymb,mathtools}
\usepackage[scr=rsfs,cal=euler]{mathalfa}
\usepackage{xfrac}
\usepackage[unicode,hidelinks,bookmarks=false]{hyperref}
\newcommand{\ie}{i.e.\ }
\newcommand{\cf}{cf.\ }
\newcommand{\resp}{respectively\ }
\newcommand{\Subsection}{\subsection}
\providecommand{\nobreakdash}{\nobreak\hbox{-}}
\setlength{\emergencystretch}{2em}
\multlinegap0pt
\usepackage{bm}
\usepackage{bbm}
\usepackage{dsfont}
\usepackage{xspace}
\usepackage{cancel}
\usepackage{mathtools}
\usepackage[shortlabels]{enumitem}
\usepackage{amsthm}
\makeatletter
\@ifundefined{theorem}{\newtheorem{theorem}{Theorem}}{}
\makeatother
\newcommand{\T}{{\mathbf{T}}}
\newcommand{\dfc}[1]{(I_{K_{\Theta_#1}}\!\!-S_{\Theta_#1}S_{\Theta_#1}^*)}
\newcommand{\hsp}{\mathcal{H}}
\newcommand{\oL }{\mathcal{L}(\hsp) }
\title[Characterizations of Matrix Valued Asymmetric Truncated Toep]{Characterizations of Matrix Valued Asymmetric Truncated Toeplitz Operators}
\author{Rewayat Khan, Yagoub Ameur,, Jamroz Khan}
\date{Source: arXiv:2103.15719v3 (CC BY 4.0)}
\begin{document}
\maketitle

\noindent\emph{Source and licence.} Characterizations of Matrix Valued Asymmetric Truncated Toeplitz Operators. Version arXiv:2103.15719v3 (CC BY 4.0), distributed under the Creative Commons Attribution 4.0 International licence, as recorded in the arXiv licence metadata for this exact version and shown on the abstract page https://arxiv.org/abs/2103.15719v3. Full rights notice: licenses/arxiv-2103.15719-CC-BY-4.0.txt.\par\smallskip

\noindent\textit{Editorial scope.} Counted unit: the Theorem whose source label is \texttt{T1} in the article (source0.tex lines 684--692, source label \texttt{T1}), delivered together with the article's own complete original proof of it (source0.tex lines 693--742, one \texttt{proof} environment of 50 lines). No mathematics is written, repaired or re-derived: statement and proof are the article's own text line for line. Required context retained uncounted: tcara (lines 472--477); cara (lines 474--476). Every definition, display and label of those lines is retained unchanged. Omitted from this package and not redistributed: 1--471, 477--683, 743--875. Nothing inside a retained excerpt is omitted or altered: every retained body is delivered exactly as the article's lines are written, so no omission receipt of any kind accompanies this package. Citations: the retained excerpts carry no citation command at all, so no bibliography entry of the article is transcribed and no reference is redistributed. Reference adaptation: none was needed and none was made; every reference inside the delivered unit resolves inside it, so no unresolved reference or citation is produced. Printed slips kept literal: no printed word, symbol, hypothesis or number of the counted statement or of its proof was altered. Layout and class: the article's own class is \texttt{amsart} with packages including epstopdf, natbib, xcolor, amsmath, amssymb, amsthm, amsfonts, arydshln, enumerate, easybmat; the wrapper is the lane's pinned \texttt{amsart} editorial wrapper at 10pt on A4, which restates the theorem-like environments the retained text needs and adds the article's own macro definitions that the retained text needs (\texttt{\textbackslash T}, \texttt{\textbackslash dfc}, \texttt{\textbackslash hsp}, \texttt{\textbackslash oL}), with extra wrapper packages (bm, bbm, dsfont, xspace, cancel, mathtools, enumitem). The delivered numbering is the unit's own. Assets: no retained span contains a figure environment, a graphics-inclusion command, a PDF-page inclusion, a table or a TikZ picture; no image file is copied into this package, the package declares no asset dependency and no image file is redistributed with it. Licence: version arXiv:2103.15719v3 of this article is distributed under the Creative Commons Attribution 4.0 International licence, as recorded in the arXiv licence metadata for this exact version and shown on the abstract page https://arxiv.org/abs/2103.15719v3. A full rights notice is delivered at licenses/arxiv-2103.15719-CC-BY-4.0.txt. Characters: every retained excerpt is pure ASCII, so the delivered engine, XeLaTeX with its default Latin Modern text font, reports no missing character.
\begin{theorem}\label{tcara}
	Let $ {\Theta_1},{\Theta_2} \in H^{\infty}(\oL)$ be two pure inner functions and let $A:K_{\Theta_{1}}\rightarrow  K_{\Theta_{2}}$ be a bounded linear operator. Then $A$ belongs to $\mathcal{MT}(\Theta_{1}, \Theta_{2})$  if and only if there exist bounded linear operators \mbox{$B_1:\mathcal{D}_{\Theta_{1}}\to K_{\Theta_{2}}$} and $B_2: \mathcal{D}_{\Theta_{2}}\to K_{\Theta_{1}}$, such that
	\begin{equation}\label{cara}
	A-S_{\Theta_2}AS_{\Theta_1}^*=B_1\dfc{1}+\dfc{2}B_2^*.
	\end{equation}
\end{theorem}

\begin{theorem}
Let $\Theta_1, \Theta_2\in H^\infty(\oL)$ be two pure inner functions and let $A:K_{\Theta_1}\to K_{\Theta_2}$ be a bounded linear operator. Then $A$ belongs to $\mathcal{MT}(\Theta_1, \Theta_2)$ if and only if it has one (and all) of the following properties:
\begin{enumerate}[(a)]
\item $\langle AS^*f, S^*g\rangle_{L^2(\hsp)}=\langle A f,g\rangle_{L^2(\hsp)}$ for all $f\in K_{\Theta_1}$, $g\in K_{\Theta_2}$ such that $f\perp\mathcal{D}_{\Theta_1}$, $g\perp\mathcal{D}_{\Theta_2}$;\label{T1}
    \item $\langle AS^*f, g\rangle_{L^2(\hsp)}=\langle Af,Sg\rangle_{L^2(\hsp)}$ for all $f\in K_{\Theta_1}$, $g\in K_{\Theta_2}$ such that $f\perp\mathcal{D}_{\Theta_1}$, $g\perp\widetilde{\mathcal{D}}_{\Theta_2}$;\label{T2}
    \item $\langle ASf, Sg\rangle_{L^2(\hsp)}=\langle Af,g\rangle_{L^2(\hsp)}$ for all $f\in K_{\Theta_1}$, $g\in K_{\Theta_2}$ such that $f\perp\widetilde{\mathcal{D}}_{\Theta_1}$, $g\perp\widetilde{\mathcal{D}}_{\Theta_2}$;\label{T3}
    \item $\langle ASf, g\rangle_{L^2(\hsp)}=\langle Af,S^*g\rangle_{L^2(\hsp)}$ for all $f\in K_{\Theta_1}$, $g\in K_{\Theta_2}$ such that $f\perp\widetilde{\mathcal{D}}_{\Theta_1}$, $g\perp\mathcal{D}_{\Theta_2}$;\label{T4}
\end{enumerate}
\end{theorem}

\begin{proof}
\eqref{T1} If $A\in \mathcal{MT}(\Theta_1, \Theta_2)$, then by Theorem \ref{tcara},
$$A-S_{\Theta_2}AS_{\Theta_1}^*=B_1D_{\Theta_1}+D_{\Theta_2}B_2^*$$
for some bounded linear operators $B_1: \mathcal{D}_{\Theta_1}\to K_{\Theta_2}$ and $B_2: \mathcal{D}_{\Theta_2}\to K_{\Theta_1}$. It follows that for all $f\in K_{\Theta_1}$, $f\in K_{\Theta_2}$ such that $f\perp\mathcal{D}_{\Theta_1}$, $g\perp\mathcal{D}_{\Theta_2}$, we have
\begin{align*}
\langle AS^*f, S^*g\rangle_{L^2(\hsp)}
&=\langle A S^*_{\Theta_1}f,S^*_{\Theta_2}g\rangle_{L^2(\hsp)}=\langle  S_{\Theta_2}AS^*_{\Theta_2}f, g\rangle_{L^2(\hsp)}\\
&=\langle Af, g\rangle_{L^2(\hsp)}-\langle B_1D_{\Theta_1}f,g\rangle_{L^2(\hsp)}-\langle D_{\Theta_2}B_2^*f,g\rangle_{L^2(\hsp)}.
\end{align*}
Since $D_{\Theta_1}f=0$ and $D_{\Theta_2}B_2^*f\in \mathcal{D}_{\Theta_2}$, we get
\begin{equation}\label{Sin1}
\langle AS^*f, S^*g\rangle_{L^2(\hsp)}=\langle A f,g\rangle_{L^2(\hsp)}
\end{equation}
On the other hand, if \eqref{Sin1} holds for all $f\in K_{\Theta_1}$, $g\in K_{\Theta_2}$ such that $f\perp \mathcal{D}_{\Theta_1}$, $g\perp \mathcal{D}_{\Theta_2}$, we have
$$\langle (A-S_{\Theta_2}AS_{\Theta_1}^*)f,g\rangle_{L^2(\hsp)}=\langle Af,g\rangle_{L^2(\hsp)}- \langle AS^*f,S^*g\rangle_{L^2(\hsp)}=0.$$
This means that the operator $\T_A=A-S_{\Theta_2}AS_{\Theta_1}^*$ maps $\mathcal{D}^{\perp}_{\Theta_1}$ into $\mathcal{D}_{\Theta_2}$, or in other words,
\begin{equation}\label{Sin2}
(I_{K_{\Theta_2}}-P_{\mathcal{D}_{\Theta_2}})\T_A(I_{K_{\Theta_1}}-P_{\mathcal{D}_{\Theta_1}})=0,
\end{equation}
where $P_{\mathcal{D}_{\Theta_i}}$ is the orthogonal projection from $K_{\Theta_i}$ onto $\mathcal{D}_{\Theta_i}$, $i=1,2$.
Recall now that
$$\text{Range}P_{\mathcal{D}_{\Theta_i}}=\mathcal{D}_{\Theta_i}=\text{Range}D_{\Theta_i},~~~ i=1,2,$$
and so there exist bounded linear operators $R_i: K_{\Theta_i}\to K_{\Theta_i}$, $i=1,2,$ such that
$$P_{\mathcal{D}_{\Theta_i}}=D_{\Theta_i}R_i=R_i^*D_{\Theta_i}, ~~~ i=1,2$$
(the second equality follows from the fact that $P_{\mathcal{D}_{\Theta_i}}^*=P_{\mathcal{D}_{\Theta_i}}$). Together with \eqref{Sin2} this gives
\begin{align*}
A-S_{\Theta_2}AS_{\Theta_1}^*
&=\T_A=P_{\mathcal{D}_{\Theta_2}}\T_A+\T_AP_{\mathcal{D}_{\Theta_2}}-P_{\mathcal{D}_{\Theta_2}}\T_AP_{\mathcal{D}_{\Theta_1}}\\
&=D_{\Theta_2}R_2\T_A+(I_{K_{\Theta_2}}-P_{\mathcal{D}_{\Theta_2}})\T_AR_1^*D_{\Theta_1}
\end{align*}
and so $A$ satisfies \eqref{cara} with
$$B_1=(I_{K_{\Theta_2}}-P_{\mathcal{D}_{\Theta_2}})
\T_AR^*_{1|\mathcal{D}_{\Theta_1}}:
\mathcal{D}_{\Theta_1}\to K_{\Theta_2}$$
and
$$B_2=(P_{\mathcal{D}_{\Theta_2}}R_2\T_A)^*=\T_A^*R^*_{2|\mathcal{D}_{\Theta_2}}:\mathcal{D}_{\Theta_2}\to K_{\Theta_1}.$$
By Theorem \ref{tcara}, $A\in \mathcal{MT}(\Theta_1, \Theta_2)$.

%Here we show that $\eqref{T2}$ is equivalent to \eqref{T1}. Assume first that $A$ satisfies \eqref{T1} and take $f\in K_{\Theta_1}$, $g\in K_{\Theta_2}$ such that $f\perp \mathcal{D}_{\Theta_1}$, $g\in \widetilde{\mathcal{D}}_{\Theta_2}$.
%Clearly, $g=S^*Sg$ and $Sg\in K_{\Theta_2}$, $Sg\perp\mathcal{D}_{\Theta_2}$. Hence
%$$\langle AS^*f, g\rangle_{L^2(\hsp)}=\langle AS^*f, S^*Sg\rangle_{L^2(\hsp)}=\langle Af, Sg\rangle_{L^2(\hsp)}$$
%and $A$ satisfies \eqref{T2}.
%
%Similarly, if $A$ satisfies \eqref{T2}, then for each $f\in K_{\Theta_1}$, $g\in K_{\Theta_2}$ such that
%$f\perp\mathcal{D}_{\Theta_1}$, $g\in \mathcal{D}_{\Theta_2}$ we have
%$$\langle AS^*f, S^*g\rangle_{L^2(\hsp)}=\langle Af, SS^*g\rangle_{L^2(\hsp)}=\langle Af, g\rangle_{L^2(\hsp)}$$
%since here $S^*f\perp\widetilde{\mathcal{D}}_{\Theta_2}$ and $S^*Sg=g$.

The proof of \eqref {T2},\eqref{T3} and \eqref{T4} is analogous to the proof of \eqref{T1}.
\end{proof}

\end{document}
